\documentclass[11pt]{article}

\usepackage[margin=1in]{geometry}
\usepackage[T1]{fontenc}
\usepackage{lmodern}
\usepackage{microtype}
\usepackage{amsmath,amssymb,amsthm,mathtools}
\usepackage{booktabs,longtable,array,tabularx}
\usepackage{enumitem}
\usepackage{xcolor}
\usepackage{hyperref}
\usepackage[nameinlink,noabbrev]{cleveref}
\usepackage{listings}
\usepackage{fancyvrb}
\usepackage{url}

\hypersetup{
  colorlinks=true,
  linkcolor=blue!45!black,
  citecolor=green!35!black,
  urlcolor=blue!55!black,
  pdftitle={All-Depth Target-One Barrier Exclusion for the Negative Collatz Shadow},
  pdfauthor={Principia Math Collatz Project}
}

\setlist{nosep,leftmargin=2.1em}
\setlength{\parindent}{0pt}
\setlength{\parskip}{0.55em}
\allowdisplaybreaks

\newtheorem{theorem}{Theorem}[section]
\newtheorem{lemma}[theorem]{Lemma}
\newtheorem{proposition}[theorem]{Proposition}
\newtheorem{corollary}[theorem]{Corollary}
\newtheorem{conjecture}[theorem]{Conjecture}
\newtheorem{problem}[theorem]{Problem}
\theoremstyle{definition}
\newtheorem{definition}[theorem]{Definition}
\newtheorem{certificate}[theorem]{Certificate}
\theoremstyle{remark}
\newtheorem{remark}[theorem]{Remark}

\newcommand{\N}{\mathbb N}
\newcommand{\Z}{\mathbb Z}
\newcommand{\Q}{\mathbb Q}
\newcommand{\R}{\mathbb R}
\newcommand{\Up}{U_{+}}
\newcommand{\Um}{U_{-}}
\newcommand{\Sp}{S_{+}}
\newcommand{\Sm}{S_{-}}
\newcommand{\vTwo}{v_2}
\newcommand{\one}{\mathbf 1}
\newcommand{\dd}{\,\mathrm d}
\newcommand{\eps}{\varepsilon}

\lstdefinelanguage{Lean}{
  morekeywords={import,open,namespace,end,section,def,abbrev,structure,inductive,
  theorem,lemma,proposition,corollary,axiom,constant,noncomputable,example,
  where,by,exact,apply,intro,have,show,from,let,if,then,else,match,with,
  forall,exists,fun,True,False,Prop,Nat,List,Finset,Set,Rat},
  sensitive=true,
  morecomment=[l]{--},
  morecomment=[s]{/-}{-/},
  morestring=[b]"
}
\lstset{
  basicstyle=\ttfamily\footnotesize,
  columns=fullflexible,
  breaklines=true,
  frame=single,
  rulecolor=\color{black!25},
  showstringspaces=false,
  tabsize=2
}

\title{\textbf{All-Depth Target-One Barrier Exclusion\\for the Negative Collatz Shadow}\\[0.4em]
\large A rigorous audit of the theorem, its certificates, the current proof architectures, and the remaining synchronization problem}
\author{Principia Math Collatz Project\\Technical audit prepared for Anton Shakov}
\date{17 August 2026}

\begin{document}
\maketitle

\begin{center}
\fbox{\begin{minipage}{0.92\textwidth}
\textbf{Status.}  The principal result of this document is an all-depth,
computer-assisted theorem about a constrained language for the accelerated
$3x-1$ map.  The proof has an analytic part and a finite exact certificate.
It is \emph{not} a proof of the Collatz conjecture.  The remaining debts are
identified precisely in \cref{sec:frontier}.
\end{minipage}}
\end{center}

\begin{abstract}
For a positive valuation word $a=(a_0,\ldots,a_{r-1})$, define
$A_k=\sum_{i<k}a_i$ and $C_0=0$, $C_{k+1}=3C_k+2^{A_k}$.
We prove that
\[
  3^r\mid C_r+2^{A_r},\qquad
  2^{A_k}<3^k\quad(1\le k\le r)
\]
forces $a_0=\cdots=a_{r-1}=1$.
Equivalently, the all-ones words are the only target-$1$ terminal words that
remain strictly on the expanding side of the coefficient barrier at every
prefix for the accelerated negative map
$\Um(m)=(3m-1)/2^{\vTwo(3m-1)}$.

The proof has four ingredients.  Terminal divisibility reconstructs an exact
positive $3x-1$ trajectory ending at $1$.  Equal-time translation on parity
cylinders gives a sharp shell-injection bound for every finite simple
negative shortcut trajectory.  This yields the universal reciprocal-mass
estimate
\[
  \sum_i\frac1{3x_i-1}<10.789.
\]
The product formula along a barrier trajectory then bounds its initial seed
by $3^{11}$ times its terminal target.  A complete scan of the odd seeds below
$200000$, implemented independently in C++ and Python, leaves only the seed
$1$.  The numerical shell sum through level $1000$, its analytic tail, every
source file, every retained output, and cryptographic hashes are specified.

We also audit the surrounding Collatz project.  We separate proved affine,
locality, coalescence, inverse-fiber, and record-compression theorems from
finite evidence, conditional completion criteria, and refuted or insufficient
routes.  The target-$1$ terminal debt is closed; nonperiodic synchronization
and positive-cycle exclusion are not.  The final appendices give Lean-oriented
proof skeletons for the all-depth theorem and for three possible routes from
the present frontier to the Collatz conjecture.
\end{abstract}

\tableofcontents
\newpage

\section{Scope, notation, and exact claim}\label{sec:scope}

The classical Collatz map is
\[
 T(n)=\begin{cases}
 n/2,&2\mid n,\\
 3n+1,&2\nmid n.
 \end{cases}
\]
It is convenient to work on odd positive integers with the accelerated map
\[
 \Up(n)=\frac{3n+1}{2^{\vTwo(3n+1)}}.
\]
The auxiliary negative map is
\[
 \Um(m)=\frac{3m-1}{2^{\vTwo(3m-1)}}.
\]
The negative map is not a replacement for the positive Collatz dynamics.  It
is a dual arithmetic shadow whose finite valuation words carry the same
powers of $2$ and $3$, with the affine carry entering with the opposite sign.

\begin{definition}[Valuation data]
Let $a=(a_0,\ldots,a_{r-1})\in\N_{\ge1}^r$.  Put
\[
 A_0=0,\qquad A_k=\sum_{i=0}^{k-1}a_i,
\]
and define
\[
 C_0=0,\qquad C_{k+1}=3C_k+2^{A_k}.
\]
Thus
\[
 C_k=\sum_{i=0}^{k-1}3^{k-1-i}2^{A_i}.
\]
The word is called \emph{barrier-admissible} when
\[
 2^{A_k}<3^k\qquad(1\le k\le r),                  \tag{1.1}
\]
and \emph{target-$y$ terminal}, for an odd positive integer $y$, when
\[
 3^r\mid C_r+2^{A_r}y.                              \tag{1.2}
\]
\end{definition}

The principal theorem is the following.

\begin{theorem}[All-depth target-$1$ barrier exclusion]\label{thm:main}
Let $a=(a_0,\ldots,a_{r-1})\in\N_{\ge1}^r$.  If
\[
 3^r\mid C_r+2^{A_r}
 \quad\text{and}\quad
 2^{A_k}<3^k\quad(1\le k\le r),
\]
then
\[
 a_0=a_1=\cdots=a_{r-1}=1.
\]
\end{theorem}

The qualifier ``all-depth'' refers to the word length $r$: no upper bound on
$r$ occurs in the statement.  The theorem is nevertheless computer-assisted,
because its analytic argument compresses every possible length to a finite
seed range and an exact finite scan settles that range.

\section{Position relative to the literature and provenance}\label{sec:literature}

The strongest unconditional global results remain far from pointwise
totality.  Tao proved that the minimum of the Collatz orbit is below every
function tending to infinity for logarithmic-density almost all starting
values~\cite{Tao2022}.  Barina's current published computation verifies
convergence for every starting value below $2^{71}$~\cite{Barina2025}.
Krasikov and Lagarias proved a power-law lower bound for the number of
starting values known to reach $1$~\cite{KrasikovLagarias2003}.  For cycles,
Simons and de Weger developed strong Diophantine and computational bounds
on the number of local minima~\cite{SimonsDeWeger2005}, and Hercher proved
that a nontrivial positive Collatz cycle would have at least $92$ local
minima~\cite{Hercher2023}.  This is the strongest peer-reviewed local-minimum
bound located in the present audit.  A June 2026 independent preprint claims
an exact computer-assisted extension excluding $m$-cycles with $m\le95$,
conditional on the Barina verification bound and the Simons--de Weger
inputs~\cite{Wang2026}; that certificate has not been independently audited
here and is not used below.

The present theorem has a different scope.  It classifies a constrained
language in the negative $3x-1$ shadow.  It neither verifies additional
ordinary positive seeds nor excludes a positive Collatz cycle.  In
particular, it is not comparable by implication with the computational bound
$2^{71}$ or Tao's almost-all theorem.

The ingredients have recognizable ancestors.

\begin{enumerate}[label=\textup{(\roman*)}]
\item Affine word formulae, parity cylinders, and $2$-adic symbolic coding are
part of the classical $3x+1$ literature; see Lagarias' survey
\cite{Lagarias1985}, Wirsching's monograph~\cite{Wirsching1998}, and the
Bernstein--Lagarias conjugacy~\cite{BernsteinLagarias1996}.
\item Orbit sparsity and reciprocal-sum arguments are closely related to the
work of Garcia and Tal on generalized $3n+1$ orbits~\cite{GarciaTal1999}.
\item The mixed binary--ternary rewrite-system viewpoint comes from
Yolcu, Aaronson, and Heule~\cite{YolcuAaronsonHeule2023}.
\item The target-shadow formulation, terminal suffix locality, coalescing
minima, and the equal-time shell injection were developed in the present
Principia Math Collatz project.  The project direction and work ledger were
maintained by Anton Shakov; derivations, audits, and executable certificates
were developed with GPT-assisted mathematical exploration and then rechecked
symbolically and computationally for this report.
\end{enumerate}

A targeted search of the standard surveys, the generalized Syracuse
literature, the rewrite-system literature, and recent work on paradoxical or
extremal words~\cite{RozierTerracol2025,FernandezIbanez2026,Kramer2026}
did not locate a published theorem equivalent to \cref{thm:main}.  This is a
search report, not a priority claim.  A formal priority claim would require a
broader expert literature review and external peer verification.

\section{Affine word calculus and terminal reconstruction}\label{sec:affine}

\subsection{Positive and negative affine identities}

\begin{proposition}[Affine identities]\label{prop:affine}
Suppose $a=(a_0,\ldots,a_{r-1})$ is the exact valuation word of a positive
accelerated trajectory $n_0,\ldots,n_r$ under $\Up$.  Then
\[
 2^{A_k}n_k=3^k n_0+C_k\qquad(0\le k\le r).          \tag{3.1}
\]
If it is the exact valuation word of a positive accelerated trajectory
$m_0,\ldots,m_r$ under $\Um$, then
\[
 2^{A_k}m_k=3^k m_0-C_k\qquad(0\le k\le r),          \tag{3.2}
\]
or equivalently
\[
 3^k m_0=2^{A_k}m_k+C_k.                              \tag{3.3}
\]
\end{proposition}

\begin{proof}
For the positive map,
$2^{a_k}n_{k+1}=3n_k+1$.  Assuming \textup{(3.1)} at time $k$ and multiplying
this recurrence by $2^{A_k}$ gives
\[
 2^{A_{k+1}}n_{k+1}
 =3(3^k n_0+C_k)+2^{A_k}
 =3^{k+1}n_0+C_{k+1}.
\]
The negative identity follows identically from
$2^{a_k}m_{k+1}=3m_k-1$.
\end{proof}

If words $u$ and $v$ are concatenated, their carries satisfy
\[
 C(uv)=3^{|v|}C(u)+2^{A(u)}C(v).                     \tag{3.4}
\]
This follows either from the defining sum or by composing the affine maps.

\subsection{Terminal suffix locality}

\begin{proposition}[Full terminal residue locality]\label{prop:locality}
Let a positive $\Up$ trajectory have endpoint $X$, and let
$v=(a_{j-r},\ldots,a_{j-1})$ be its last $r$ valuation exponents.  Put
$B=A(v)$ and $C_v=C(v)$.  Then, for every integer $c$,
\[
 X\equiv c\pmod{3^r}
 \quad\Longleftrightarrow\quad
 C_v-2^B c\equiv0\pmod{3^r}.                         \tag{3.5}
\]
In particular,
\[
 3^r\mid X+1
 \quad\Longleftrightarrow\quad
 3^r\mid C_v+2^B.                                    \tag{3.6}
\]
\end{proposition}

\begin{proof}
If $Y$ is the positive odd state at the start of the suffix, then
\[
 2^B X=3^rY+C_v.
\]
Reduction modulo $3^r$ and invertibility of $2^B$ modulo $3^r$ prove the
claim.
\end{proof}

\subsection{Terminal divisibility reconstructs an exact trajectory}

\begin{theorem}[Targeted terminal duality]\label{thm:duality}
Let $a=(a_0,\ldots,a_{r-1})\in\N_{\ge1}^r$ and let $y$ be a positive odd
integer.  The following are equivalent.
\begin{enumerate}[label=\textup{(\roman*)}]
\item $3^r\mid C_r+2^{A_r}y$.
\item There are positive odd integers $m_0,\ldots,m_r=y$ satisfying
\[
 3m_i-1=2^{a_i}m_{i+1}\qquad(0\le i<r).              \tag{3.7}
\]
\end{enumerate}
Every exponent in \textup{(3.7)} is then the exact $2$-adic valuation of
$3m_i-1$.
\end{theorem}

\begin{proof}
For $0\le i\le r$, let $D_i$ denote
\[
 D_i=C(a_i,\ldots,a_{r-1})
      +2^{a_i+\cdots+a_{r-1}}y,
\]
with $D_r=y$.  Splitting off the first letter of a suffix gives
\[
 D_i=3^{r-i-1}+2^{a_i}D_{i+1}.                       \tag{3.8}
\]
Assume $3^{r-i}\mid D_i$.  Reducing \textup{(3.8)} modulo
$3^{r-i-1}$ shows that $3^{r-i-1}\mid D_{i+1}$.  Starting from (i), this
propagates divisibility to every suffix.  Define
\[
 m_i=\frac{D_i}{3^{r-i}}.
\]
Then \textup{(3.8)} becomes
\[
 3m_i=1+2^{a_i}m_{i+1}.
\]
Backward induction from the positive odd integer $m_r=y$ shows that each
$m_i$ is a positive odd integer: the numerator $1+2^{a_i}m_{i+1}$ is odd,
and divisibility by $3$ has already been proved.  Consequently
$3m_i-1=2^{a_i}m_{i+1}$ with $m_{i+1}$ odd, so the valuation is exactly
$a_i$.  Conversely, iterating \textup{(3.7)} gives
$3^rm_0=C_r+2^{A_r}y$.
\end{proof}

This theorem is the first decisive reduction: a formal congruence on a word
is not merely necessary.  It produces an ordinary positive integer
trajectory in the negative map.

\section{Equal-time parity cylinders}\label{sec:equaltime}

Define the negative shortcut map on nonnegative integers by
\[
 \Sm(x)=\begin{cases}
 x/2,&x\equiv0\pmod2,\\
 (3x-1)/2,&x\equiv1\pmod2.
 \end{cases}                                          \tag{4.1}
\]
For $k\ge0$, let $s_k(x)$ be the number of odd inputs among
$x,\Sm(x),\ldots,\Sm^{k-1}(x)$.

\begin{lemma}[Parity cylinder]\label{lem:parity-cylinder}
Congruence modulo $2^k$ determines the first $k$ parity bits of a shortcut
trajectory.  Moreover, every binary word of length $k$ occurs for exactly one
residue class modulo $2^k$.
\end{lemma}

\begin{proof}
If $x\equiv y\pmod{2^k}$, then $x$ and $y$ have the same parity.  After one
common branch, their images are congruent modulo $2^{k-1}$.  Induction gives
the first assertion.

For injectivity, suppose two residues have the same parity word, of weight
$o$.  Following the common affine branch word gives
\[
 \Sm^k(x)-\Sm^k(y)=\frac{3^o}{2^k}(x-y).
\]
The left side is an integer and $3^o$ is odd, so $2^k\mid x-y$.  The two
residues are equal.  There are $2^k$ residues and $2^k$ binary words, hence
the map is bijective.
\end{proof}

\begin{lemma}[Equal-time translation]\label{lem:equal-time}
For nonnegative integers $x,q,k$,
\[
 \Sm^k(x+q2^k)=\Sm^k(x)+q3^{s_k(x)}.                 \tag{4.2}
\]
\end{lemma}

\begin{proof}
The two initial values are congruent modulo $2^k$, so by
\cref{lem:parity-cylinder} they follow the same first $k$ branches.  Their
initial difference is $q2^k$.  Each step divides the difference by $2$, and
each odd branch additionally multiplies it by $3$.  After $k$ steps the
difference is $q3^{s_k(x)}$.
\end{proof}

For a fixed parity word of length $k$ and weight $o$, the corresponding
affine formula has the form
\[
 2^k\Sm^k(u)=3^o u-D,\qquad D\ge0.                   \tag{4.3}
\]
Therefore, whenever $0\le u<2^k$,
\[
 0\le\Sm^k(u)<3^o.                                   \tag{4.4}
\]
The nonnegativity follows because the shortcut map preserves
$\N\cup\{0\}$; the strict upper bound follows from $u<2^k$ and $D\ge0$.

\section{Finite-orbit shell sparsity}\label{sec:shell}

Let $x_0,x_1,\ldots,x_L$ be a positive $\Sm$ trajectory with no repeated
state.  For $k\ge0$, put
\[
 N_k=\#\{i:2^k\le x_i<2^{k+1}\}.                    \tag{5.1}
\]
Define
\[
 B_k=\min_{0\le t\le k}
 \left(
   \sum_{o=t+1}^k\binom{k}{o}
   +\sum_{o=0}^t3^o
 \right).                                             \tag{5.2}
\]

\begin{theorem}[Shell sparsity]\label{thm:shell}
For every $k\ge0$,
\[
 N_k\le\min\{2^k,B_k+k\}.                            \tag{5.3}
\]
\end{theorem}

\begin{proof}
The bound $N_k\le2^k$ is immediate because the shell
$[2^k,2^{k+1})$ contains exactly $2^k$ integers.

Fix $t\in\{0,\ldots,k\}$.  Temporarily discard the final $k$ positions of
the trajectory.  Every remaining shell point has a unique representation
\[
 x_i=2^k+u_i,
 \qquad0\le u_i<2^k.                                  \tag{5.4}
\]
By \cref{lem:parity-cylinder}, exactly $\binom{k}{o}$ residues $u$ have a
first-$k$ parity word of weight $o$.  Hence the number of retained shell
points whose weight exceeds $t$ is at most
\[
 \sum_{o=t+1}^k\binom{k}{o}.                          \tag{5.5}
\]

Now fix $o\le t$.  For a retained shell point of weight $o$, define its
\emph{base endpoint} to be $\Sm^k(u_i)$.  By \textup{(4.4)}, the base endpoint
is one of the $3^o$ integers in $[0,3^o)$.  We claim that, for fixed $o$, no
two retained shell points have the same base endpoint.  Indeed, if
$x_i=2^k+u_i$ and $x_j=2^k+u_j$ had the same weight and the same base
endpoint, \cref{lem:equal-time} with $q=1$ would give
\[
 x_{i+k}=\Sm^k(u_i)+3^o
        =\Sm^k(u_j)+3^o=x_{j+k}.
\]
Both future positions exist because the last $k$ positions were discarded.
This contradicts nonrepetition.  Thus the number of retained points of
weight $o$ is at most $3^o$.  Summing over $o\le t$, adding \textup{(5.5)},
and restoring at most $k$ discarded positions gives
\[
 N_k\le
 \sum_{o=t+1}^k\binom{k}{o}+\sum_{o=0}^t3^o+k.
\]
Minimizing over $t$ proves \textup{(5.3)}.
\end{proof}

The role of equal time is essential.  It turns a collision in a finite set
of base endpoints into an actual repeated trajectory state.  Counting parity
words alone would not provide this injection.

\section{A universal reciprocal-mass constant}\label{sec:mass}

Every point in the $k$th shell contributes at most
$(3\cdot2^k-1)^{-1}$.  Therefore \cref{thm:shell} gives
\[
 \sum_{i=0}^{L}\frac1{3x_i-1}
 \le
 \sum_{k=0}^{\infty}
 \frac{\min\{2^k,B_k+k\}}{3\cdot2^k-1}.              \tag{6.1}
\]
We now prove that the right side is below $10.789$.

\subsection{The analytic tail}

\begin{lemma}[Chernoff form for the binomial tail]\label{lem:chernoff}
Let $1/2<p<1$ and let $t=\lfloor pk\rfloor$.  Then
\[
 \sum_{o=t+1}^k\binom{k}{o}
 \le\left(p^{-p}(1-p)^{-(1-p)}\right)^k.             \tag{6.2}
\]
\end{lemma}

\begin{proof}
For $z>1$ and every $o>pk$, one has $z^o>z^{pk}$.  Hence
\[
 \sum_{o>pk}\binom{k}{o}
 \le z^{-pk}\sum_{o=0}^k\binom{k}{o}z^o
 =z^{-pk}(1+z)^k.
\]
The right side is minimized at $z=p/(1-p)$, which yields \textup{(6.2)}.
\end{proof}

Take $p=61/100$ and $t=\lfloor61k/100\rfloor$.  The exact integer
inequalities
\[
 3^{61}<2^{97},                                        \tag{6.3}
\]
\[
 100^{100}<2^{97}61^{61}39^{39}                       \tag{6.4}
\]
are equivalent to
\[
 3^{61/100}<2^{97/100},
 \qquad
 p^{-p}(1-p)^{-(1-p)}<2^{97/100}.
\]
By \cref{lem:chernoff},
\[
 \sum_{o=t+1}^k\binom{k}{o}<2^{97k/100}.              \tag{6.5}
\]
Also,
\[
 \sum_{o=0}^t3^o=\frac{3^{t+1}-1}{2}
 <\frac32 3^t
 <\frac32 2^{97k/100}.                                \tag{6.6}
\]
Consequently
\[
 B_k<\frac52 2^{97k/100}.                              \tag{6.7}
\]
For $k\ge1001$,
\[
 \frac{B_k+k}{3\cdot2^k-1}
 <\frac54 2^{-3k/100}+\frac{k}{2^{k+1}}.              \tag{6.8}
\]
Grouping the first series into blocks of $100$ terms gives
\[
 \begin{aligned}
 \frac54\sum_{k=1001}^{\infty}2^{-3k/100}
 &<\frac54\sum_{j=0}^{\infty}100\,2^{-30-3j}\\
 &=\frac{1000}{7\cdot2^{30}}.
 \end{aligned}                                        \tag{6.9}
\]
The standard differentiated geometric series gives
\[
 \sum_{k=1001}^{\infty}\frac{k}{2^{k+1}}
 =\frac{1002}{2^{1001}}.                               \tag{6.10}
\]
Thus the exact tail bound is
\[
 \sum_{k=1001}^{\infty}
 \frac{\min\{2^k,B_k+k\}}{3\cdot2^k-1}
 <\frac{1000}{7\cdot2^{30}}+\frac{1002}{2^{1001}}
 <1.331\times10^{-7}.                                  \tag{6.11}
\]

\subsection{The finite shell certificate}

For $0\le k\le1000$, define
\[
 n_k=\min\{2^k,B_k+k\},
 \qquad
 u_k=\left\lceil
 10^{15}\frac{n_k}{3\cdot2^k-1}
 \right\rceil.                                        \tag{6.12}
\]
The retained CSV certificate contains $(k,B_k,t_k,n_k,u_k)$ for every
$k$, where $t_k$ is the first minimizing threshold in \textup{(5.2)}.
An independent verifier rebuilds Pascal's triangle, recomputes every
threshold, and checks every row.

\begin{certificate}[Finite mass certificate]\label{cert:mass}
The exact verified integer identity is
\[
 \sum_{k=0}^{1000}u_k
 =10{,}788{,}901{,}977{,}820{,}654.                   \tag{6.13}
\]
Therefore
\[
 \sum_{k=0}^{1000}\frac{n_k}{3\cdot2^k-1}
 \le10.788901977820654.                               \tag{6.14}
\]
Combining \textup{(6.11)} and \textup{(6.14)} gives the certified rational upper
bound
\[
 10.788902110866736\ldots<10.789.                     \tag{6.15}
\]
\end{certificate}

No floating-point comparison enters the proof.  Decimal values in
\textup{(6.14)}--\textup{(6.15)} are displays of exact rational inequalities.  The
proof-relevant data are the integers in \textup{(6.3)}, \textup{(6.4)},
\textup{(6.13)}, and the rational tail in \textup{(6.11)}.

\begin{theorem}[Universal reciprocal mass]\label{thm:mass}
Every finite positive simple trajectory
$x_0,\ldots,x_L$ of the negative shortcut map satisfies
\[
 \sum_{i=0}^{L}\frac1{3x_i-1}<10.789.                 \tag{6.16}
\]
\end{theorem}

\begin{proof}
Combine \textup{(6.1)}, \cref{cert:mass}, and \textup{(6.11)}.
\end{proof}

\section{Fixed-target compression and proof of the main theorem}\label{sec:compression}

\begin{theorem}[Fixed-target seed compression]\label{thm:compression}
Let
\[
 m_0\longmapsto m_1\longmapsto\cdots\longmapsto m_r=y
\]
be an exact trajectory of $\Um$ on positive odd integers whose expansion to shortcut steps has no repeated state, with valuation
word $a$ satisfying
\[
 2^{A_j}<3^j\qquad(1\le j\le r).
\]
Then
\[
 m_0<e^{10.789}y<3^{11}y=177147y.                    \tag{7.1}
\]
\end{theorem}

\begin{proof}
Each accelerated step satisfies
\[
 \frac{m_{j+1}}{m_j}
 =\frac{3m_j-1}{2^{a_j}m_j}
 =\frac3{2^{a_j}}\left(1-\frac1{3m_j}\right).
\]
Multiplication gives
\[
 \frac{y}{m_0}
 =\frac{3^r}{2^{A_r}}
 \prod_{j=0}^{r-1}\left(1-\frac1{3m_j}\right).       \tag{7.2}
\]
The full barrier gives $3^r/2^{A_r}>1$, hence
\[
 \log\frac{m_0}{y}
 <-\sum_{j=0}^{r-1}\log\left(1-\frac1{3m_j}\right). \tag{7.3}
\]
For $0<u<1$,
\[
 -\log(1-u)=\int_0^u\frac{\dd t}{1-t}
 \le\frac{u}{1-u}.
\]
With $u=1/(3m_j)$,
\[
 -\log\left(1-\frac1{3m_j}\right)
 \le\frac1{3m_j-1}.                                  \tag{7.4}
\]
Expand the accelerated trajectory into shortcut steps.  Before the terminal
state this shortcut trajectory is simple, and the accelerated odd states are
a subset of its states.  Thus \cref{thm:mass} gives
\[
 \log\frac{m_0}{y}<10.789.
\]
Finally $10.789<11$ and $e<3$ (for example, from the exponential series), so $e^{10.789}<3^{11}$.
\end{proof}

The final finite input is the following exact computation.

\begin{certificate}[Target-$1$ seed scan]\label{cert:seed}
Every odd seed $m_0$ with $1\le m_0<200000$ was followed under exact integer
arithmetic until the first occurrence of one of the following mutually
exclusive events:
\begin{enumerate}[label=\textup{(\roman*)}]
\item the orbit reaches $1$;
\item a state repeats;
\item the strict accelerated barrier fails;
\item the explicit step cap is reached; or
\item a checked fixed-width operation would overflow.
\end{enumerate}
The C++ implementation returned
\[
 \begin{array}{lr}
 \text{reaches }1&1\\
 \text{repeated-state cycles}&2\\
 \text{barrier failures}&99997\\
 \text{step-cap survivors}&0\\
 \text{overflows}&0.
 \end{array}
\]
The unique terminal hit was $(m_0,r)=(1,0)$.  A separate Python
implementation on the shortcut map, using arbitrary-precision integers,
returned the same partition.
\end{certificate}

\begin{proof}[Proof of \cref{thm:main}]
By \cref{thm:duality}, terminal divisibility reconstructs a positive exact
trajectory
\[
 m_0\longmapsto\cdots\longmapsto m_r=1
\]
under $\Um$.  Let $\tau$ be the first index with $m_\tau=1$.  If
$\tau>0$, then the trajectory up to $m_\tau$ has no repeated state: a repeat
in a deterministic map would trap the orbit in a cycle and prevent a later
first hit of $1$.  The prefix barrier remains valid.  Applying
\cref{thm:compression} to this truncated path gives
\[
 1<m_0<177147<200000.
\]
This contradicts \cref{cert:seed}, which has no such seed.  Hence
$\tau=0$ and $m_0=1$.

Now $3\cdot1-1=2$.  Since the next state is odd, the exact valuation is
$a_0=1$ and $m_1=1$.  Repeating this argument proves
$a_i=1$ and $m_i=1$ for every $i<r$.
\end{proof}

\section{Certificate specification and reproducibility}\label{sec:certificates}

Only the certificates in this section are logical dependencies of
\cref{thm:main}.  Earlier barrier-tree, order, and excursion computations are
useful evidence but are not used in the proof.

\subsection{Retained files and commands}

The companion archive contains:
\begin{center}
\begin{tabularx}{\textwidth}{@{}lX@{}}
\toprule
File & Function\\
\midrule
\texttt{mass\_certificate.py} & Computes every $B_k$ for $k\le1000$, the upward-rounded rational sum, and the exact analytic-tail inequalities.\\
\texttt{mass\_shells\_k0\_1000.csv} & The 1001-row finite shell certificate.\\
\texttt{verify\_mass\_csv.py} & Independent Pascal-row verifier for every CSV field and the total.\\
\texttt{accelerated\_seed\_certificate.cpp} & Exact accelerated $\Um$ scan using checked \texttt{unsigned \_\_int128} multiplication and arbitrary-precision barrier powers.\\
\texttt{shortcut\_seed\_certificate.py} & Independent shortcut scan using Python arbitrary-precision integers.\\
\texttt{*.out} & Retained standard output for each run.\\
\texttt{ENVIRONMENT.txt} & Python, compiler, and operating-system identifiers.\\
\texttt{SHA256SUMS} & Cryptographic digest of every retained proof artifact.\\
\bottomrule
\end{tabularx}
\end{center}

The reproduction commands are
\begin{Verbatim}[fontsize=\small,frame=single]
python mass_certificate.py --K 1000 --csv mass_shells_k0_1000.csv
python verify_mass_csv.py
python shortcut_seed_certificate.py --cap 200000 --step-cap 100000
g++ -O2 -std=c++20 -Wall -Wextra -Werror \
  accelerated_seed_certificate.cpp -o accelerated_seed_certificate
./accelerated_seed_certificate 200000 100000
sha256sum -c SHA256SUMS
\end{Verbatim}

\subsection{Exact retained outputs}

The mass computation reports
\begin{Verbatim}[fontsize=\small,frame=single]
scaled_upward_sum=10788901977820654
check_3pow61_lt_2pow97=True
check_entropy_integer_inequality=True
tail_upper_decimal=1.330460820879255e-07
certified_total_upper_decimal=10.788902110866736
certified_total_lt_10_789=True
\end{Verbatim}
The independent CSV verifier reports
\begin{Verbatim}[fontsize=\small,frame=single]
rows_verified=1001
scaled_upward_sum=10788901977820654
certificate_passed=True
\end{Verbatim}
The shortcut and accelerated seed scans report, respectively,
\begin{Verbatim}[fontsize=\small,frame=single]
status_counts={'barrier_failure': 99997, 'cycle': 2, 'hit_1': 1}
terminal_hits=[(1, 0)]
longest_resolution_steps=132
longest_resolution_seed=149345
certificate_passed=True
\end{Verbatim}
\begin{Verbatim}[fontsize=\small,frame=single]
hit_1=1
cycle=2
barrier_failure=99997
step_cap=0
overflow=0
longest_resolution_steps=83
longest_resolution_seed=149345
terminal_hits= (1,0)
certificate_passed=True
\end{Verbatim}
The different longest-resolution lengths are expected: one program counts
accelerated odd steps, while the other counts every shortcut step.

\subsection{SHA--256 manifest}

\begin{center}
\scriptsize
\begin{tabularx}{\textwidth}{@{}lX@{}}
\toprule
File & SHA--256\\
\midrule
\texttt{mass\_certificate.py} & \texttt{8d83628d918c9828207be42beb26633b770722d4561c2db7e4e2afcb0856c7fa}\\
\texttt{verify\_mass\_csv.py} & \texttt{76426c4f5400e98fcc31af15d119a54cbd5268d836291b2816429a841082759f}\\
\texttt{shortcut\_seed\_certificate.py} & \texttt{615e0badd5b6ce48834b9541ab2e3d010526a5cdba5d24565cea0f284d06a45a}\\
\texttt{accelerated\_seed\_certificate.cpp} & \texttt{4b9bfa5bca84fdd064f1fee0cb85e8ce74985fa06c6bb1dee98fe67ec91c4620}\\
\texttt{mass\_shells\_k0\_1000.csv} & \texttt{3997b941ae24e59e262636d51ab5baa82b6262de0cd04f6a5812978e7ebea266}\\
\texttt{mass\_certificate.out} & \texttt{8630b315ec64a84983041f428f804f48c56be27b3a122ab52ec593bd76acb21f}\\
\texttt{verify\_mass\_csv.out} & \texttt{e7ceaab0085fd8f32368b26df80ff05709f203f95608687839f47e246d91c032}\\
\texttt{shortcut\_seed\_certificate.out} & \texttt{efb67cb920fd8cc9a89e2a4410c1a8ab47dc6560d96f792709abf8b81fcb56f8}\\
\texttt{accelerated\_seed\_certificate.out} & \texttt{30f5dd9f84454be45a7840b6a3a52f0edbfeccb42c4fb8e6cf9fdd42e86a3348}\\
\texttt{ENVIRONMENT.txt} & \texttt{3df8386c359b194a6414fe70b80e4231b3cfd018f9c02924b2e7c07f356e353c}\\
\bottomrule
\end{tabularx}
\end{center}

\subsection{Trust boundary}

The finite theorem can be checked at several levels.

\begin{enumerate}[label=\textup{(\roman*)}]
\item A reader may inspect the mathematical reduction and trust the retained
executables and outputs.
\item A reader may recompile and rerun both independent seed scans and both
mass computations.
\item A formalization can replace the CSV verifier and seed loops by
\texttt{native\_decide}, while proving the analytic shell and tail lemmas in
the Lean kernel.  A detailed skeleton appears in \cref{app:lean}.
\end{enumerate}

The theorem does not rely on probabilistic sampling, floating-point
rounding, an unreported height cutoff, or an extrapolation in the word
length.

\section{What the theorem represents and what it does not}\label{sec:meaning}

\subsection{A local language theorem of genuinely infinite depth}

Before \cref{thm:main}, the project had an exact finite exclusion through
word length $10^{10}$.  That result still left the logical possibility of a
first nontrivial target-$1$ barrier word at a larger length.  The new argument
changes the quantifier structure: arbitrary word length is converted into a
uniform seed bound, and the finite scan is over seeds rather than lengths.
This is why the theorem is genuinely all-depth.

The key compression is not a generic property of symbolic languages.  It
uses three features simultaneously:
\begin{enumerate}[label=\textup{(\alph*)}]
\item exact parity-cylinder translation at a common time;
\item positivity of the negative affine carry, which bounds base endpoints;
\item the multiplicative correction product along an exact $3x-1$ path.
\end{enumerate}
Removing any of these features destroys the present proof.

\subsection{Consequence for positive Collatz suffixes}

Combining \cref{prop:locality,thm:main} gives the following exact statement.

\begin{corollary}[Rate-matched terminal suffixes]\label{cor:positive-suffix}
Let a positive accelerated Collatz block of length $r$ have valuation word
$v$, and suppose every relative prefix satisfies
$2^{A_k(v)}<3^k$.  If its endpoint $X$ satisfies
$3^r\mid X+1$, then $v=1^r$.
\end{corollary}

Thus every nontrivial barrier block has
\[
 v_3(X+1)<r.                                           \tag{9.1}
\]
The theorem classifies the exact $3$-adic target hit.  It does not prove that
a hypothetical divergent positive orbit must produce such a hit.

\subsection{Why positive cycles remain separate}

If a positive accelerated cycle of length $r$ has total exponent $A$ and
least point $x$, its full affine equation is
\[
 (2^A-3^r)x=C_r>0.                                    \tag{9.2}
\]
Hence $2^A>3^r$, opposite to the full barrier inequality in
\cref{thm:main}.  No rotation of the theorem can simply erase this sign
change.

There is also an exact obstruction to ``pumping'' a positive cycle into the
target-$1$ language.  Let a shortcut-cycle word $w$ have length $p$, odd
weight $s$, $Q=2^p$, $P=3^s$, and cycle point $x$, so
$C_w=(Q-P)x$.  For the $N$th power,
\[
 C_{w^N}=x(Q^N-P^N).
\]
Appending an all-odd word $1^L$ gives
\[
 C_{w^N1^L}+Q^N2^L
 =3^L\bigl((x+1)Q^N-xP^N\bigr).                      \tag{9.3}
\]
Consequently $w^N1^L$ is target-$1$ terminal exactly when
\[
 P^N\mid x+1.                                         \tag{9.4}
\]
The number of admissible repetitions is therefore bounded by
$v_3(x+1)/s$.  Periodic positive dynamics cannot be converted into
arbitrarily deep target-$1$ words by this attachment.

\section{Audit of the rigorous project ledger}\label{sec:ledger}

This section consolidates the recent manuscripts, proof skeletons, and
certificate reports.  The status labels are deliberately strict.

\subsection{Established exact identities and reductions}

\begin{theorem}[Two-ended gluing]\label{thm:gluing}
Let $a$ label an exact negative trajectory
$m_0\to\cdots\to m_r=y$, and let $A=A_r$.  For every positive even integer
$s$ such that $n=2^As-m_0>0$, the same word is the exact positive valuation
word from $n$ to
\[
 X=3^rs-y.                                             \tag{10.1}
\]
At intermediate times,
\[
 n_k=2^{A-A_k}3^ks-m_k.                                \tag{10.2}
\]
\end{theorem}

\begin{proof}
The formula holds at $k=0$.  If it holds at $k$, then
\[
 \begin{aligned}
 3n_k+1
 &=2^{A-A_k}3^ks\cdot3-(3m_k-1)\\
 &=2^{a_k}\bigl(2^{A-A_{k+1}}3^{k+1}s-m_{k+1}\bigr).
 \end{aligned}
\]
The bracket is $n_{k+1}$.  Since $s$ is even and $m_{k+1}$ is odd, every
$n_{k+1}$ is odd; hence the displayed power of $2$ is exact.
\end{proof}

This theorem parametrizes coalescing positive ancestors by negative shadows.
It leads to the exact min-plus inverse-tree formulation used later.

For a positive word prefix, put
\[
 q_j=\frac{3^j}{2^{A_j}},\qquad
 \beta_j=\frac{C_j}{2^{A_j}},
\]
so that
\[
 x_j=q_jn+\beta_j.                                    \tag{10.3}
\]
When $q_j<1$, define
\[
 R_j=\frac{\beta_j}{1-q_j}
 =\frac{C_j}{2^{A_j}-3^j}.                            \tag{10.4}
\]
Then
\[
 x_j-n=(1-q_j)(R_j-n).                                \tag{10.5}
\]
Thus a least-counterexample orbit that never falls below its start must have
$R_j\ge n$ at every negative-coefficient prefix.

\begin{proposition}[Block-threshold convexity]\label{prop:convexity}
Let a prefix $u$ have multiplier $t<1$ and normalized carry $y$, and append a
block $v$ with multiplier $Q<1$ and normalized carry $\gamma$.  Put
\[
 R(u)=\frac{y}{1-t},\qquad R(v)=\frac{\gamma}{1-Q}.
\]
Then
\[
 R(uv)=\lambda R(u)+(1-\lambda)R(v),
 \qquad
 \lambda=\frac{Q(1-t)}{1-Qt}\in(0,1).                \tag{10.6}
\]
\end{proposition}

\begin{proof}
The affine pair updates by
\[
 (t,y)\longmapsto(Qt,Qy+\gamma).
\]
Substitute $y=(1-t)R(u)$ and $\gamma=(1-Q)R(v)$ into
$(Qy+\gamma)/(1-Qt)$.
\end{proof}

\begin{theorem}[Stationary inverse-fiber extinction]\label{thm:SFE}
For an odd positive integer $n$, define
\[
 M_r(n)=\min\{z>0:z\text{ odd and }\Up^r(z)=\Up^r(n)\}. \tag{10.7}
\]
The Collatz conjecture is equivalent to
\[
 \forall n>1\text{ odd}\quad\exists r\ge1
 \quad M_r(n)<n.                                      \tag{10.8}
\]
\end{theorem}

\begin{proof}
If $n$ reaches $1$ after $r$ accelerated steps, then the fixed point $1$ is
an $r$-step ancestor of the same endpoint, so $M_r(n)=1<n$.
Conversely, assume \textup{(10.8)}.  Strong induction on odd $n$ proves
convergence.  Choose $r$ with $z=M_r(n)<n$.  By induction, $z$ reaches $1$.
Since $z$ and $n$ have the same endpoint after $r$ accelerated steps, $n$
also reaches $1$.
\end{proof}

This is the narrowest exact one-statement reformulation presently available
in the project.  It absorbs both cycles and divergent orbits, but its forward
inequality is not yet proved.

\subsection{Divergent-orbit compression}

Let $x_j=\Up^j(n)$ be a nonperiodic positive odd orbit with least value $n$.
Set
\[
 h_j=\frac{x_j}{q_j}.
\]
From the affine recurrence,
\[
 h_{j+1}=h_j\left(1+\frac1{3x_j}\right).              \tag{10.9}
\]
Apart from a bounded initial contribution, the distinct odd endpoints lie in
the two residue classes $1$ and $5$ modulo $6$.  Ordering the smallest
possible values in those classes gives
\[
 \sum_{i<j}\frac1{x_i}\le\frac13\log j+O_n(1).       \tag{10.10}
\]
Therefore
\[
 \log\frac{h_j}{n}
 \le\frac13\sum_{i<j}\frac1{x_i}
 \le\frac19\log j+O_n(1),                            \tag{10.11}
\]
and hence
\[
 h_j\le K_nj^{1/9},\qquad
 q_j\ge c_nj^{-1/9}.                                  \tag{10.12}
\]
Equivalently,
\[
 A_j\log2-j\log3\le\frac19\log j+O_n(1).             \tag{10.13}
\]
At a strict lower record of $q_j$, one has $q_j\le1$ and therefore
\[
 x_j=q_jh_j\le K_nj^{1/9}.                             \tag{10.14}
\]
Together with
$C_j\equiv2^{A_j}x_j\pmod{3^j}$, this reduces record endpoints and exponent
surplus to polynomially many possibilities.  The valuation word is uniquely
recoverable from $(j,A_j,C_j)$ by repeatedly using
\[
 a_0=\vTwo(C_j-3^{j-1})                                \tag{10.15}
\]
and stripping the first letter.  This is the project's
\emph{record-endpoint compression theorem}.

A second exact input uses Garcia--Tal orbit sparsity: a nonperiodic orbit that
does not reach the trivial cycle tends to infinity and has a convergent
reciprocal sum.  The product corresponding to \textup{(10.9)} then converges,
so $q_j\to\infty$.  After shifting to a global minimum of $q$, the relative
coefficient sequence satisfies $q_j\ge1$ for every $j$ and tends to infinity.
In the inverse-carry coordinates of the parity conjugacy, ordinary positive
integers impose eventual finite support.  The divergence problem is thereby
reduced to excluding the coexistence of unbounded critical surplus and an
eventually zero inverse-carry stream.

\subsection{Exact rewrite macrosemantics}

The mixed binary--ternary rewrite system of
Yolcu--Aaronson--Heule simulates Collatz~\cite{YolcuAaronsonHeule2023}.
The project proved an elementary ordinal phase rank below $\omega^3$ for the
ten nonexceptional rules and proved that their canonical normal form is the
binary numeral for the odd part of the represented integer.  Consequently,
one exceptional rewrite followed by complete normalization is exactly one
accelerated $\Up$ step, independently of normalization order.  This removes
syntactic nondeterminism but also exposes the limit of the approach: the
quotient macrograph is the accelerated Collatz graph itself.  A termination
proof for the complete rewrite system still requires a genuine arithmetic
descent theorem.

\subsection{Finite certificates retained as evidence but not used in
\cref{thm:main}}

\begin{longtable}{@{}p{0.27\textwidth}p{0.67\textwidth}@{}}
\toprule
Certificate & Exact finite conclusion\\
\midrule
\endfirsthead
\toprule
Certificate & Exact finite conclusion\\
\midrule
\endhead
Previous target-$1$ word scan & The all-ones word was the only target-$1$ barrier word for every length $r\le10^{10}$.  Superseded by \cref{thm:main}.\\
Barrier-order certificate & At depth $21$, $39{,}993{,}895$ barrier words were grouped by endpoint residue.  There were no equal-total-exponent collisions and no violations of the conjectured coalescing order; the minimum positive margin was $32$.\\
Compressed charged-excursion scan & For every $3$-free quotient $q\le10^9$, except the known nondescending cycle states $q=2,8$, the first descent found within the cap obeyed the tested inequality $L\le2(q-Q^L(q))$; no violations occurred.\\
Negative $d$-bound scan & For all odd seeds $m_0\le99{,}999{,}999$ whose negative orbit terminates at $1$, the tested inequality $2d\le m_0-1$ held, where $d$ counts exponents exceeding $1$.  This is evidence, not an all-seed proof.\\
Barrier trees and suffix locality & Exact finite trees through the reported depths, and hundreds of thousands of suffix-locality comparisons, produced no mismatch.  These computations motivated the exact locality theorem but are no longer needed once it is proved symbolically.\\
\bottomrule
\end{longtable}

\subsection{Exact no-go results and routes that were not completed}

The project also made progress by proving that several attractive shortcuts
cannot carry the missing conclusion in their proposed form.

\begin{longtable}{@{}p{0.24\textwidth}p{0.70\textwidth}@{}}
\toprule
Route & Audited outcome\\
\midrule
\endfirsthead
\toprule
Route & Audited outcome\\
\midrule
\endhead
Arbitrary target recurrence & False as a direct shortcut: the negative cycle $5\leftrightarrow7$ makes $(1,2)^N$ a nontrivial barrier word terminal at target $5$ for every $N$.  Target recurrence must exclude periodic negative shadows.\\
Ordinary compactness & Two sufficiently close terminal prefixes for one target force the intervening block to be a negative-map cycle.  Thus bounded-gap or subexponential recurrence tends to return exactly the wrong periodic objects.\\
Cycle pumping into target $1$ & Ruled out by \textup{(9.4)}: only finitely many repetitions can satisfy target-$1$ divisibility.\\
Finite automata alone & Every finite parity word is realizable by a residue class modulo a power of $2$.  Any successful automaton must retain unbounded arithmetic carry or height information.\\
Local thermodynamic potential & A continuous finite-memory ``integer-lift potential'' cannot distinguish ordinary positive integers from general $2$-adic parity sequences; the positive-integer lift is nonlocal.\\
Cobham-style dual automaticity & The natural transducers act on value digits rather than the time index, and periodic cycles would already satisfy eventual automaticity.\\
Abstract $C^*$-simplicity & The abstract algebra is already the relevant Cuntz algebra independently of Collatz; the meaningful concrete commutant reduces to functions constant on orbit components.\\
Finite modular/K-theoretic approximations & Natural modular graphs and the full $2$-adic boundary lose the positive-integer lifting condition and take the same coarse value for systems with known nontrivial cycles.\\
Canonical height & Natural polarization, Northcott, and product-formula requirements are incompatible with the powers-of-two basin and the local slopes; no coercive point-height of the proposed form exists.\\
Pure entropy or random-word drift & Average negative drift does not imply a uniform theorem on the thin set of words realized by one fixed exceptional integer.  The missing step is a pointwise synchronization or inverse-fiber extinction statement.\\
\bottomrule
\end{longtable}

These negative results are useful because they prevent the proof skeleton from
silently replacing an ordinary-integer theorem by a statement about random or
$2$-adic words.

\section{The current proof frontier}\label{sec:frontier}

The all-depth theorem closes one previously open component:
\[
 \boxed{\text{nontrivial target-$1$ terminal/barrier words do not exist}.}
\]
It does not close the two global debts.

\subsection{Architecture A: terminal exclusion plus synchronization plus
cycle exclusion}

A sufficient proof architecture is
\[
 \begin{gathered}
 \text{positive-cycle exclusion}\\
 +\quad\text{nonperiodic synchronization}\\
 +\quad\text{all-depth target-$1$ exclusion}\\
 \Longrightarrow\quad\text{Collatz}.
 \end{gathered}                                        \tag{11.1}
\]
The third line is \cref{thm:main}.  The other two are open.

A granular synchronization theorem strong enough for \textup{(11.1)} would say:
for a globally least unbounded counterexample $n$, there exist a block start
$j$ and a length $r$ such that, with relative sums
$B_k=A_{j+k}-A_j$,
\[
 2^{B_k}<3^k\quad(1\le k\le r),                       \tag{11.2}
\]
\[
 3^r\mid x_{j+r}+1,                                   \tag{11.3}
\]
and the block word is not $1^r$.  By suffix locality, \textup{(11.3)} makes the
word target-$1$ terminal; \cref{thm:main} then contradicts nontriviality.

Each clause matters.
\begin{enumerate}[label=\textup{(\roman*)}]
\item Barrier-block abundance alone gives \textup{(11.2)} for arbitrarily long
blocks on a divergent orbit, but gives no $3$-adic endpoint condition.
\item Generic $3$-adic recurrence can give endpoint congruences, but close
repetitions can collapse to periodic negative shadows.
\item An all-ones block satisfies both the barrier and the target condition,
so synchronization must force nontriviality or provide an additional descent
from long all-ones blocks.
\item The synchronization must concern one fixed ordinary integer.  An
average over seeds, or a statement for random parity words, is insufficient.
\end{enumerate}

The cycle debt is equally exact: exclude every primitive word $a$ and odd
$x>1$ satisfying
\[
 (2^{A_r}-3^r)x=C_r                                   \tag{11.4}
\]
with all cyclic valuation constraints.  Hercher's $92$-local-minimum theorem is the strongest peer-reviewed
structural lower bound verified in this audit.  The more recent conditional
preprint claim through $m=95$ is recorded in \cref{sec:literature} but is not
used.  No uniform classification of every possible number of minima is known.

\subsection{Architecture B: stationary inverse-fiber extinction}

By \cref{thm:SFE}, it is enough to prove the single statement
\[
 \boxed{
 \forall n>1\text{ odd}\quad\exists r\ge1\quad M_r(n)<n.
 }                                                       \tag{11.5}
\]
This formulation automatically handles both cycles and divergence.  It asks
for a smaller positive ancestor that coalesces with the actual orbit after a
finite time.  Through \cref{thm:gluing}, every candidate ancestor can be
written as the cost of a negative shadow.  The exact missing theorem is that
the positive minimum in this target-shadow min-plus tree cannot remain
stationary forever.

This is the logically narrowest completion target, but it is also equivalent
to Collatz.  Its value is organizational: it tells every proposed method
exactly which finite witness it must construct.

\subsection{Architecture C: first crossing and high-surplus carry innovation}

The stationary-inverse-fiber manuscript isolates a more structured pair of
lemmas.

\begin{conjecture}[First-crossing descent]\label{conj:FCD}
Whenever the coefficient $q_j=3^j/2^{A_j}$ first crosses below $1$ along a
positive ordinary orbit, the associated threshold/carry data produce either
a smaller coalescing ancestor or a quantitatively controlled successor state
in the normalized barrier regime.
\end{conjecture}

\begin{conjecture}[High-surplus carry innovation]\label{conj:HSCI}
No infinite ordinary parity word can simultaneously have unbounded critical
surplus in the normalized barrier regime and an eventually zero inverse-carry
stream.  Equivalently, sufficiently high surplus forces a new nonzero binary
lift digit or a smaller coalescing ancestor.
\end{conjecture}

Together with the exact coefficient-record dichotomy, these two statements
imply Collatz.  They separate the finite first crossing from the infinite
normalized tail.  Neither has been proved.

\subsection{A practical hierarchy of the remaining tasks}

The next work should be judged against the following checklist.

\begin{enumerate}[label=\textbf{D\arabic*.}]
\item \textbf{Formalize the all-depth theorem.}  Prove parity cylinders,
shell injection, the Chernoff tail, and fixed-target compression in Lean;
verify the finite data with \texttt{native\_decide}.  This does not advance
the frontier mathematically, but removes implementation trust.
\item \textbf{Decide the cycle route.}  Either prove a uniform local-minimum
classification extending Hercher to every fixed number, or prove a different
uniform exclusion of solutions to \textup{(11.4)}.
\item \textbf{Upgrade barrier-block abundance to arithmetic endpoint
synchronization.}  Control $v_3(x_{j+r}+1)$ at a rate proportional to $r$,
not merely along a sparse average or a fixed modulus.
\item \textbf{Exclude the all-ones escape.}  If rate-matched target hits can
be all ones, prove that arbitrarily long such blocks force a smaller
coalescing ancestor, periodicity, or eventual entry into the trivial basin.
\item \textbf{Alternatively prove inverse-fiber saving directly.}  Show that
an infinite positive slope-geodesic cannot remain within a bounded exponent
loss of the minimum at every depth.
\item \textbf{Preserve ordinary-integer lifting.}  Every compactness,
entropy, spectral, or automata argument must retain both the eventually zero
binary start lift and the terminal ternary carry.  Losing either one returns
false $2$-adic countermodels.
\end{enumerate}

\section{Discussion and assessment}\label{sec:assessment}

\subsection{Why this is genuine progress}

The theorem removes an unbounded word-length quantifier from a central
subproblem.  Previous computations could only move the cutoff.  The new
shell argument proves that a terminal barrier path cannot begin arbitrarily
high, regardless of its length.  That is a structural theorem, not a larger
enumeration.

It also exposes a reusable principle.  For every fixed odd target $y$, every
simple target-$y$ barrier path begins below $177147y$.  Hence the entire
nonperiodic target language for a fixed $y$ is finitely classifiable; only
negative cycles can generate arbitrarily long periodic target shadows.  This
sharpens the synchronization problem: a successful recurrence theorem must
separate aperiodic target hits from periodic negative components.

\subsection{Why it is not yet close enough to declare Collatz solved}

The positive map does not automatically produce the negative target-$1$
words excluded by the theorem.  The exact obstruction is cross-prime
synchronization.  A positive ordinary start fixes an eventually constant
binary lift.  A target-$1$ terminal suffix fixes a long ternary congruence at
the endpoint.  The same valuation word must satisfy both constraints while
remaining above the least positive orbit value.  No existing theorem cited
here supplies that simultaneous pointwise statement.

Cycles are not a technical afterthought.  Their full-period coefficient has
the opposite sign, and periodic negative shadows are genuine counterexamples
to naive arbitrary-target recurrence.  Any complete paper must either absorb
cycles into a single inverse-fiber theorem or prove a separate uniform cycle
exclusion.

\subsection{Most promising exact targets}

The three best-defined next targets are:
\begin{enumerate}[label=\textup{(\arabic*)}]
\item stationary inverse-fiber extinction \textup{(11.5)};
\item nontrivial rate-matched synchronization
\textup{(11.2)}--\textup{(11.3)}, together with cycle exclusion;
\item high-surplus carry innovation in the normalized divergent regime.
\end{enumerate}
The first is the cleanest equivalence, the second uses the new all-depth
theorem most directly, and the third preserves the strongest quantitative
compression already proved for divergent orbits.

\appendix

\section{Dependency graph for the all-depth theorem}\label{app:dependency}

\begin{Verbatim}[fontsize=\small,frame=single]
affine carry recursion
        |
        +--> terminal divisibility reconstructs exact U_- path
        |
parity-cylinder bijection
        |
        +--> equal-time translation
                    |
                    +--> shell injection N_k <= min(2^k,B_k+k)
                                      |
                    exact k<=1000 sum + analytic k>=1001 tail
                                      |
                    universal reciprocal mass < 10.789
                                      |
negative product formula + full barrier
        |
        +--> fixed-target seed compression m_0 < 3^11 y
                                      |
                    y=1 + exact odd seed scan below 200000
                                      |
                    only m_0=1
                                      |
                    all exponents are 1
\end{Verbatim}

No theorem below the seed scan depends on any prior finite word-depth bound.
No synchronization or positive-cycle hypothesis enters the proof.

\section{Pseudocode for the exact certificates}\label{app:pseudocode}

\subsection{Mass certificate}

\begin{lstlisting}[language=Python]
SCALE = 10**15
scaled_sum = 0
for k in range(0, 1001):
    binom = [C(k,o) for o in range(k+1)]
    Bk = min(
        sum(binom[o] for o in range(t+1,k+1))
        + sum(3**o for o in range(0,t+1))
        for t in range(k+1)
    )
    n = min(2**k, Bk + k)
    u = ceil(SCALE*n / (3*2**k - 1))   # exact integer ceiling
    scaled_sum += u
assert scaled_sum == 10788901977820654
assert 3**61 < 2**97
assert 100**100 < 2**97 * 61**61 * 39**39
assert scaled_sum/SCALE + 1000/(7*2**30) + 1002/2**1001 < 10.789
\end{lstlisting}

\subsection{Accelerated seed certificate}

\begin{lstlisting}[language=C++]
for (odd seed = 1; seed < 200000; seed += 2) {
    m = seed; A = 0; seen = empty;
    while (true) {
        if (m == 1) classify(seed, HIT_ONE);
        if (m in seen) classify(seed, CYCLE);
        insert m into seen;
        z = checked_exact_value(3*m - 1);
        a = v2(z);
        m = z >> a;
        A += a; r += 1;
        if (2^A >= 3^r) classify(seed, BARRIER_FAILURE);
    }
}
assert counts == {HIT_ONE:1, CYCLE:2, BARRIER_FAILURE:99997};
assert terminal_hits == {(1,0)};
\end{lstlisting}

The actual C++ certificate stores $2^A$ and $3^r$ in arbitrary-precision
integers and computes $3m-1$ in a checked 128-bit temporary.  The independent
Python shortcut scan does not share this implementation.

\section{Lean-oriented proof and completion skeletons}\label{app:lean}

The full pseudocode file \texttt{collatz\_proof\_skeleton.lean} accompanies
this paper.  It is deliberately a blueprint rather than a claim of completed
formalization.  The most important declarations are reproduced below.

\lstinputlisting[language=Lean]{Collatz_Proof_Skeleton.lean}

\begin{thebibliography}{99}

\bibitem{Barina2025}
D. Barina,
\newblock Improved verification limit for the convergence of the Collatz conjecture,
\newblock \emph{J. Supercomput.} \textbf{81} (2025), Article 810,
\newblock doi:10.1007/s11227-025-07337-0.

\bibitem{BernsteinLagarias1996}
D. J. Bernstein and J. C. Lagarias,
\newblock The $3x+1$ conjugacy map,
\newblock \emph{Canad. J. Math.} \textbf{48} (1996), 1154--1169,
\newblock doi:10.4153/CJM-1996-060-x.

\bibitem{FernandezIbanez2026}
C. Fern\'andez and S. Ib\'a\~nez,
\newblock Christoffel words as extremal structures in Collatz dynamics,
\newblock arXiv:2607.24844, 2026.

\bibitem{GarciaTal1999}
M. V. P. Garcia and F. A. Tal,
\newblock A note on the generalized $3n+1$ problem,
\newblock \emph{Acta Arith.} \textbf{90} (1999), 245--250.

\bibitem{Hercher2023}
C. Hercher,
\newblock There are no Collatz $m$-cycles with $m\le91$,
\newblock \emph{J. Integer Sequences} \textbf{26} (2023), Article 23.3.5.

\bibitem{Kramer2026}
O. Kramer,
\newblock Adaptive search in Collatz exponent-code space via $2$-adic and $3$-adic constraints,
\newblock arXiv:2607.10041, 2026.

\bibitem{KrasikovLagarias2003}
I. Krasikov and J. C. Lagarias,
\newblock Bounds for the $3x+1$ problem using difference inequalities,
\newblock \emph{Acta Arith.} \textbf{109} (2003), 237--258,
\newblock doi:10.4064/aa109-3-4.

\bibitem{Lagarias1985}
J. C. Lagarias,
\newblock The $3x+1$ problem and its generalizations,
\newblock \emph{Amer. Math. Monthly} \textbf{92} (1985), 3--23.

\bibitem{RozierTerracol2025}
O. Rozier and C. Terracol,
\newblock Paradoxical behavior in Collatz sequences,
\newblock \emph{Discrete Math.} \textbf{349} (2026), 115167,
\newblock doi:10.1016/j.disc.2026.115167.

\bibitem{SimonsDeWeger2005}
J. Simons and B. de Weger,
\newblock Theoretical and computational bounds for $m$-cycles of the $3n+1$ problem,
\newblock \emph{Acta Arith.} \textbf{117} (2005), 51--70,
\newblock doi:10.4064/aa117-1-3.

\bibitem{Tao2022}
T. Tao,
\newblock Almost all orbits of the Collatz map attain almost bounded values,
\newblock \emph{Forum Math. Pi} \textbf{10} (2022), e12,
\newblock doi:10.1017/fmp.2022.8.


\bibitem{Wang2026}
X. Wang,
\newblock Non-existence of Collatz $m$-cycles for $m\le95$,
\newblock independent preprint, June 2026,
\newblock doi:10.5281/zenodo.21017296.

\bibitem{Wirsching1998}
G. J. Wirsching,
\newblock \emph{The Dynamical System Generated by the $3n+1$ Function},
\newblock Lecture Notes in Mathematics 1681, Springer, 1998.

\bibitem{YolcuAaronsonHeule2023}
E. Yolcu, S. Aaronson, and M. J. H. Heule,
\newblock An automated approach to the Collatz conjecture,
\newblock \emph{J. Autom. Reason.} \textbf{67} (2023), Article 15,
\newblock doi:10.1007/s10817-022-09658-8.

\end{thebibliography}

\end{document}
